\subsection{Thông tin tài liệu và lịch sử phiên bản}
\label{sec:document-control}

\begin{guide}
\guidepart{Nội dung}
\begin{compactitem}
    \item Bảng thông tin tài liệu: tên tài liệu, phiên bản hiện tại, trạng thái (Draft, Review, Final), ngày cập nhật, nhóm tác giả, người duyệt.
    \item Bảng lịch sử phiên bản: phiên bản, ngày, người thực hiện, mô tả thay đổi. Mỗi lần nộp (release \texttt{phase-*} trên GitHub) tương ứng với một dòng trong bảng.
    \item Mục đích, phạm vi và đối tượng đọc của tài liệu, cách tổ chức các phần còn lại theo chuẩn IEEE 1016~\cite{ieee1016}.
    \item Bảng thuật ngữ và từ viết tắt dùng trong tài liệu (SDD, MVP, IA, USP, a11y, \ldots).
\end{compactitem}
\guidepart{Câu hỏi cần trả lời}
\begin{compactitem}
    \item Tài liệu mô tả thiết kế của phiên bản nào, giai đoạn nào của ứng dụng?
    \item Ai là người đọc chính (giảng viên, nhóm phát triển, tester) và họ dùng tài liệu này để làm gì?
    \item Quy ước đánh số phiên bản là gì (ví dụ 0.x là bản nháp, 1.0 là bản nộp)? Mỗi phiên bản thay đổi những phần nào?
\end{compactitem}
\guidepart{Hình/bảng nên có}
Bảng thông tin tài liệu, bảng lịch sử phiên bản, bảng thuật ngữ.
\end{guide}
